\documentclass{amsart}

\input{../preamble_general.tex}

\author{Jalen Chrysos}
\title{Rooted Trees \& Cayley's Formula}

\begin{document}
	\maketitle 

%	A subgraph $T$ of $G$ is called a \textit{spanning tree} if $T$ includes every vertex and is a tree. For each graph $G$, let $t(G)$ be the number of spanning trees of $G$.\\
%	
%	\noindent \textbf{Problem}: Show if $G$ is connected then $t(G)\geq 1$.\\
%	
%	\noindent \textbf{Problem}: Show that if a tree $T$ is a subgraph of $G$ (not necessarily a spanning tree), then $T$ can be extended to a spanning tree of $G$.\\
%	
%	\noindent \textbf{Problem}: Show that if $C_n$ is a subgraph of $G$, $t(G)\geq n$.\\
%	
%	\noindent \textbf{Problem}: In this problem we will show Cayley's Formula, which states that $t(K_n)=n^{n-2}$. This proof, due to Jim Pitman, calculates $t(K_n)$ by first calculating the number of \textit{rooted} spanning trees, which turns out to be much easier.
%	\begin{enumerate}[(a)]
%		\item A \textit{rooted tree} is a tree along with orientations of each edge such that every vertex has out-degree at most 1.
%	\end{enumerate}


	A \textit{rooted tree} is a tree $T$ with a distinguished vertex $r$ called the ``root.'' \\
	
	\noindent \textbf{Problem 1}: Given a tree $T$, the choice of a root $r$ induces an orientation on the edges of $T$: 
	\begin{enumerate}[(a)]
		\item Show that for every edge $e\in T$, with endpoints $x_1,x_2$, there is a path in $T$ either of the form
		$$(x_1,x_2,\dots,r) \;\; \text{or} \;\; (x_2,x_1,\dots,r),$$
		but not both. Thus we can orient $e$ uniquely so that it is part of a directed path ending at $r$.
		\item Show that under the orientation induced by $r$ described in (a), every vertex has out-degree exactly 1 except for $r$ itself.
		\item Show conversely that if $T$ is a directed tree such that every vertex has out-degree $\leq 1$, then this orientation is the same as one induced by some choice of a root in $T$.\\
	\end{enumerate} 
	
	\noindent \textbf{Problem 2}: In this problem we will derive Cayley's Formula, which states that the graph $K_n$ has exactly $n^{n-2}$ spanning trees. This proof was first described by Jim Pitman.
	\begin{enumerate}[(a)]
		\item Call a nested sequence of subgraphs 
		$$A_1\subset A_2\subset \cdots \subset A_{n-1}$$
		an \textit{assembly} if each $A_j$ is a directed subgraph of $K_n$ with exactly $j$ edges. Show that for any rooted spanning tree $T\subset K_n$, there is an assembly such that $A_{n-1}=T$.
		\item For each rooted spanning tree $T$, how many assemblies are there with $T=A_{n-1}$?
		\item Compute the number of assemblies such that $A_{n-1}$ is a rooted tree. (\textit{Hint: consider adding a directed edge to $A_j$: how many places can the head go? How many places can the tail go?})
		\item Using (b) and (c), compute the number of rooted spanning trees of $K_n$.
		\item Finally, using (d), compute the number of spanning trees of $K_n$.\\
	\end{enumerate}
	
\end{document}